\documentclass[dvipdfmx,aspectratio=169]{beamer}
\usepackage{amsmath}
\usepackage{booktabs}
\usepackage{pxjahyper}
\usepackage{tikz}
\usetikzlibrary{arrows.meta,positioning,backgrounds,fit}

\usetheme{Madrid}
\usecolortheme{default}
\setbeamertemplate{navigation symbols}{}

\title[ABC128 C]{AtCoder ABC128 -- C\\\small Switches の解説}
\author{}
\date[ABC128]{AtCoder Beginner Contest 128 / 300 点}

% よく使う色
\definecolor{ok}{RGB}{0,140,70}
\definecolor{ng}{RGB}{200,30,30}
\definecolor{onc}{RGB}{240,180,0}

\begin{document}

\begin{frame}
  \titlepage
\end{frame}

%====================================================================
\begin{frame}{問題の概要}
  \begin{itemize}
    \item スイッチが $N$ 個（各 ON/OFF），電球が $M$ 個ある．
    \item 電球 $i$ は $k_i$ 個のスイッチ $s_{i,1},\ldots,s_{i,k_i}$ につながっている．
    \item \textbf{電球 $i$ が点灯する条件}：つながっているスイッチのうち
      \textbf{ON の個数} $\equiv p_i \pmod 2$．
      \begin{itemize}
        \item $p_i = 0$ なら「ON が偶数本」，$p_i = 1$ なら「ON が奇数本」．
      \end{itemize}
    \item \textbf{求めるもの}：$M$ 個すべてが点灯する ON/OFF の組み合わせの\textbf{総数}．
  \end{itemize}

  \vspace{0.6em}
  \textbf{制約}: $1 \le N, M \le 10$,\quad $1 \le k_i \le N$,\quad $p_i \in \{0,1\}$

  \vspace{0.4em}
  \begin{block}{この制約が答えを教えている}
    $N \le 10$ $\Rightarrow$ 組み合わせは高々 $2^{10} = 1024$ 通り．\textbf{全部試してよい．}
  \end{block}
\end{frame}

%====================================================================
\begin{frame}{まず制約を読む -- 全探索が許されるか見積もる}
  \begin{columns}
    \begin{column}{0.54\textwidth}
      スイッチ 1 個につき ON/OFF の 2 通りなので，全体で
      \[
        2^N \le 2^{10} = 1024 \ \text{通り}
      \]
      1 通りあたりの判定コストは
      \[
        \textstyle\sum_i k_i \le M\cdot N \le 100
      \]
      合計しても $1024 \times 100 \approx 10^5$ 回．
      \vspace{0.5em}

      \begin{alertblock}{典型パターン}
        「$N \le 20$ 程度」$+$「各要素が 2 択」\\
        $\Rightarrow$ \textbf{bit 全探索を疑う}
      \end{alertblock}
    \end{column}
    \begin{column}{0.44\textwidth}
      \centering
      \begin{tikzpicture}[scale=0.8]
        \draw[-{Latex}] (0,0) -- (0,4.3) node[above]{$2^N$};
        \draw[-{Latex}] (0,0) -- (5.2,0) node[right]{$N$};
        \foreach \n/\h in {2/0.2,4/0.4,6/0.8,8/1.6,10/3.2}{
          \fill[ok!60] (\n*0.45-0.15,0) rectangle (\n*0.45+0.15,\h);
          \node[font=\tiny] at (\n*0.45,-0.25) {\n};
        }
        \node[font=\scriptsize,anchor=west] at (4.6,3.2) {$1024$};
        \draw[ng,dashed] (0,4.0) -- (5.0,4.0)
          node[right,font=\scriptsize,ng,anchor=south east] at (5.0,4.0) {実行時間の壁 (遠い)};
      \end{tikzpicture}

      \vspace{0.3em}
      {\scriptsize $N=10$ でも $1024$．壁までまだ何桁も余裕がある．}
    \end{column}
  \end{columns}
\end{frame}

%====================================================================
\begin{frame}{アイデア(1) -- ON/OFF の並びを整数 1 個で表す}
  $N$ 個の ON/OFF は長さ $N$ の $0/1$ 列．これを整数 \texttt{bit} の 2 進表記とみなす．

  \vspace{0.5em}
  \begin{center}
    \begin{tikzpicture}[scale=1.0]
      % ビット箱
      \foreach \i/\v/\c in {0/1/onc, 1/0/white, 2/1/onc}{
        \draw[thick] (\i*1.9,0) rectangle (\i*1.9+1.2,0.9);
        \fill[\c] (\i*1.9+0.02,0.02) rectangle (\i*1.9+1.18,0.88);
        \draw[thick] (\i*1.9,0) rectangle (\i*1.9+1.2,0.9);
        \node at (\i*1.9+0.6,0.45) {\v};
      }
      \node[font=\scriptsize] at (0.6,-0.35) {第 2 bit};
      \node[font=\scriptsize] at (2.5,-0.35) {第 1 bit};
      \node[font=\scriptsize] at (4.4,-0.35) {第 0 bit};
      \node[font=\scriptsize] at (0.6,1.3) {sw2 : ON};
      \node[font=\scriptsize] at (2.5,1.3) {sw1 : OFF};
      \node[font=\scriptsize] at (4.4,1.3) {sw0 : ON};
      \node[anchor=west] at (5.9,0.45) {$\texttt{bit} = 101_{(2)} = 5$};
    \end{tikzpicture}
  \end{center}

  \vspace{0.3em}
  \begin{columns}
    \begin{column}{0.5\textwidth}
      \centering
      {\small
      \begin{tabular}{ccccc}
        \toprule
        \texttt{bit} & 2進 & sw2 & sw1 & sw0\\
        \midrule
        0 & 000 & OFF & OFF & OFF\\
        1 & 001 & OFF & OFF & \textbf{ON}\\
        2 & 010 & OFF & \textbf{ON} & OFF\\
        $\vdots$ & & & & \\
        7 & 111 & \textbf{ON} & \textbf{ON} & \textbf{ON}\\
        \bottomrule
      \end{tabular}}
    \end{column}
    \begin{column}{0.46\textwidth}
      \texttt{bit} を $0$ から $2^N-1$ まで動かせば，
      \textbf{全パターンをちょうど 1 回ずつ}列挙できる．

      \vspace{0.5em}
      「スイッチ $v$ は ON か？」は
      \[
        \texttt{bit \& (1 << v)} \ne 0
      \]
      で判定できる．
    \end{column}
  \end{columns}
\end{frame}

%====================================================================
\begin{frame}{アイデア(2) -- 1 パターンを固定して電球を判定する}
  \begin{columns}
    \begin{column}{0.46\textwidth}
      \centering
      \begin{tikzpicture}[scale=0.9,
        sw/.style={draw,thick,minimum width=0.9cm,minimum height=0.5cm,font=\small},
        bulb/.style={draw,thick,circle,minimum size=0.8cm,font=\small}]
        \node[font=\scriptsize,anchor=west] at (-0.5,2.9) {黄 $=$ ON};
        \node[sw,fill=onc] (s1) at (0,2.2) {sw1};
        \node[sw]          (s2) at (0,1.2) {sw2};
        \node[sw,fill=onc] (s3) at (0,0.2) {sw3};
        \node[bulb,fill=ok!25] (b1) at (2.6,1.9) {$b_1$};
        \node[bulb,fill=ng!20] (b2) at (2.6,0.5) {$b_2$};
        \draw (s1) -- (b1); \draw (s2) -- (b1); \draw (s3) -- (b1);
        \draw (s2) -- (b2); \draw (s3) -- (b2);
      \end{tikzpicture}

      \vspace{0.4em}
      {\scriptsize
        $b_1$: ON は 2 本 $\Rightarrow 0$，$p_1{=}0$ \textcolor{ok}{\textbf{OK}}\\
        $b_2$: ON は 1 本 $\Rightarrow 1$，$p_2{=}0$ \textcolor{ng}{\textbf{NG}}\\
        $\Rightarrow$ このパターンは失格}
    \end{column}
    \begin{column}{0.52\textwidth}
      \texttt{bit} を 1 つ固定したら，電球 $i$ ごとに
      \begin{enumerate}
        \item つながっているスイッチのうち ON の本数 \texttt{con} を数える
        \item \texttt{con \% 2 == p[i]} なら点灯，違えば\textbf{このパターンは失格}
      \end{enumerate}

      \vspace{0.5em}
      $M$ 個すべて点灯したパターンだけを数え上げれば，それが答え．

      \vspace{0.5em}
      \begin{block}{計算量}
        $O(2^N \cdot \sum k_i) \le 1024 \times 100 \approx 10^5$
      \end{block}
    \end{column}
  \end{columns}
\end{frame}

%====================================================================
\begin{frame}[fragile]{擬似コードと実装の要点}
\begin{verbatim}
res = 0
for bit in 0 .. 2^N - 1:          # 全パターン列挙
    ok = true
    for i in 0 .. M-1:
        con = (s[i] のうち bit で ON のスイッチの本数)
        if con % 2 != p[i]: ok = false
    if ok: res += 1
\end{verbatim}

  \vspace{0.2em}
  実装で効くのは 3 か所だけ:
  \begin{itemize}
    \item \verb|int a; cin >> a; --a;| \quad 入力は 1-indexed．
      \textbf{0-indexed に直してビット位置と一致}させる．
      忘れるとスイッチ $N$ が範囲外ビットを指し，常に OFF 扱いになる．
    \item \verb|for (int bit = 0; bit < (1<<N); ++bit)| \quad
      $2^N$ が全パターン数．$\le$ ではなく $<$．
    \item \verb|if (bit & (1<<v)) ++con;| \quad
      返るのは $0$ か\textbf{非 $0$}（$1$ とは限らない）．真偽値として使う．
  \end{itemize}
\end{frame}

%====================================================================
\begin{frame}{サンプル 1 を全部書き出す（答え: 1）}
  入力: $N=2$, 電球1 = \{sw1, sw2\} で $p_1=0$（偶数），電球2 = \{sw2\} で $p_2=1$（奇数）

  \vspace{0.5em}
  \centering
  \begin{tabular}{ccccccc}
    \toprule
    \texttt{bit} & sw1 & sw2 & 電球1 の ON 本数 & $\equiv p_1{=}0$? & 電球2 の ON 本数 & $\equiv p_2{=}1$?\\
    \midrule
    00 & OFF & OFF & 0 & \textcolor{ok}{OK}  & 0 & \textcolor{ng}{NG}\\
    01 & ON  & OFF & 1 & \textcolor{ng}{NG}  & 0 & \textcolor{ng}{NG}\\
    10 & OFF & ON  & 1 & \textcolor{ng}{NG}  & 1 & \textcolor{ok}{OK}\\
    \textbf{11} & \textbf{ON} & \textbf{ON} & \textbf{2} & \textcolor{ok}{\textbf{OK}} & \textbf{1} & \textcolor{ok}{\textbf{OK}}\\
    \bottomrule
  \end{tabular}

  \vspace{0.8em}
  \raggedright
  両方 OK なのは \texttt{bit}$=11$ の 1 通りだけ $\Rightarrow$ 出力 \textbf{1}．

  \vspace{0.4em}
  {\small サンプル 3 ($N=5$, 条件 2 本) の答え $8$ は，
  \textbf{どの電球にもつながっていない sw4} が自由（$\times 2$）で，
  残りの自由度も含めて $2^3 = 8$．}
\end{frame}

%====================================================================
\begin{frame}{一段深い見方 -- 実は GF(2) 上の連立一次方程式}
  スイッチ $v$ の状態を $x_v \in \{0,1\}$ と置くと，電球 $i$ の条件はそのまま
  \[
    x_{s_{i,1}} \oplus x_{s_{i,2}} \oplus \cdots \oplus x_{s_{i,k_i}} = p_i
    \qquad (\oplus = \text{XOR} = \bmod 2 \text{ の足し算})
  \]

  つまり \textbf{$\bmod 2$ の連立一次方程式の解の個数}を数える問題．線形代数から
  \[
    \text{解の個数} =
    \begin{cases}
      0 & (\text{矛盾するとき})\\[2pt]
      2^{\,N - \mathrm{rank}} & (\text{矛盾しないとき})
    \end{cases}
  \]

  \vspace{0.3em}
  \centering
  \begin{tabular}{ccccc}
    \toprule
    サンプル & $N$ & 式の本数 & rank & 解の個数\\
    \midrule
    1 & 2 & 2 & 2 & $2^{2-2} = \mathbf{1}$\\
    2 & 2 & 3 & \multicolumn{2}{c}{矛盾あり $\Rightarrow \mathbf{0}$}\\
    3 & 5 & 2 & 2 & $2^{5-2} = \mathbf{8}$\\
    \bottomrule
  \end{tabular}

  \vspace{0.6em}
  \raggedright
  {\small 本問は $N \le 10$ なので全探索で十分．\textbf{$N$ が $10^5$ 規模なら
  掃き出し法（Gauss--Jordan over GF(2)，\texttt{bitset} で高速化）}に切り替える．}
\end{frame}

%====================================================================
\begin{frame}[fragile]{まとめ}
  \begin{enumerate}
    \item \textbf{制約を先に読む．} $N \le 20$ 程度 $+$ 各要素 2 択なら bit 全探索．
    \item \textbf{bit 全探索の型}：
      \verb|for (int bit = 0; bit < (1<<N); ++bit)| と \verb|bit & (1<<v)|．
    \item \textbf{入力の 1-indexed は読み込み時に \texttt{-{}-a} で潰す．}
      番号とビット位置を一致させる．
    \item \textbf{XOR 条件の数え上げ $=$ GF(2) の連立方程式．}
      答えは $0$ か $2^{N-\mathrm{rank}}$．全探索が間に合わない規模ならこちら．
  \end{enumerate}

  \vspace{0.8em}
  \textbf{同系統の問題}
  \begin{itemize}
    \item ABC079 C -- Train Ticket（$\pm$ の全探索）
    \item ABC147 C -- HonestOrUnkind2（$2^N$ 全探索 $+$ 整合性判定．骨格がほぼ同じ）
    \item ABC167 C -- Skill Up（$2^N \times$ コスト計算）
  \end{itemize}
\end{frame}

\end{document}
